\documentclass[12pt]{article}
\usepackage[spanish]{babel}
\usepackage[T1]{fontenc}
\usepackage[utf8]{inputenc}
\usepackage{geometry}
\usepackage{booktabs}
\usepackage{array}
\usepackage{longtable}
\usepackage[hidelinks]{hyperref}
\usepackage{enumitem}

\geometry{margin=1in}
\setlist{itemsep=0.2em, topsep=0.2em}

\renewcommand{\refname}{Referencias}

\begin{document}
\begin{titlepage}
\centering
\vspace*{2cm}
{\Large UPNA - Minería de Datos\par}
\vspace{0.6cm}
\rule{0.8\textwidth}{0.6pt}\par
\vspace{0.6cm}
{\LARGE\bfseries Modelo predictor del autismo\par}
\vspace{0.6cm}
\rule{0.8\textwidth}{0.6pt}\par
\vspace{0.6cm}
{\large Eloy Urriens\\Juan Felipe Pulgarín Lopez\\Nicolás Chareca\par}
\vspace{0.6cm}
{\large 13 de mayo de 2026\par}
\vspace*{1.5cm}
\begin{abstract}
Este informe documenta el desarrollo de un modelo de Inteligencia Artificial, utilizando técnicas de minería de datos para la
predicción del autismo en personas usando el conjunto de datos
\texttt{Autism-Adult-Data.csv}. Se describe la limpieza y preparación mediante un
\texttt{Preparador}, el análisis exploratorio de variables y la definición de componentes
del pipeline (imputación, transformación, detección de outliers, estandarización
y modelos base). Se especifica la metodología de validación basada en
Pipeline + GridSearchCV con estratificación, y las métricas priorizadas
(F1-score y AUC de Precisión-Recall). El documento establece el marco
metodológico y el flujo de trabajo para el desarrollo y evaluación del modelo.
\end{abstract}
\end{titlepage}


\clearpage

\tableofcontents

\clearpage
\leftrightarrow 
\section{Introducción}
\subsection{Problema}
El objetivo del proyecto es la \textbf{predicción del autismo en personas} a partir del conjunto de datos
\texttt{Autism-Adult-Data.csv} \cite{autism_adults_kaggle,thabtah2017asd,thabtah2017review}. El trabajo se estructura en tres pasos:
\begin{itemize}
    \item Analizar el problema y las variables disponibles.
    \item Construir la mejor \emph{pipeline} posible.
    \item Reflexionar sobre la calidad de la solución.
\end{itemize}

\subsubsection*{Preparador}
Se implementa un \emph{Preparador} cuyo objetivo es ordenar y limpiar el dataset sin modificar la
información relevante. En concreto, realiza:
\begin{itemize}
    \item Renombrado de columnas y corrección de nombres erróneos.
    \item Eliminación de columnas innecesarias.
    \item Inversión del significado en un subconjunto de preguntas.
    \item Limpieza de texto en variables categóricas.
    \item Conversión de valores perdidos ``?'' a \texttt{NaN}.
    \item Tipado de variables: categóricas y numéricas.
\end{itemize}

\subsection{Análisis de variables}
Se analiza la relación entre las variables de entrada y la variable de salida.
Cuando se indica $\uparrow$ Autismo o $\downarrow$ Autismo se refiere a un aumento o
reducción de la probabilidad de pertenecer a la clase positiva.

\subsubsection*{Relación de las 10 preguntas con la variable de salida}
\begin{longtable}{p{0.06\textwidth} p{0.66\textwidth} p{0.22\textwidth}}
\toprule
\# & Pregunta & Dirección \\
\midrule
1 & I often notice small sounds when others do not. & Agree $\rightarrow$ $\uparrow$ Autismo \\
2 & I usually concentrate more on the whole picture, rather than the small details. & Agree $\rightarrow$ $\downarrow$ Autismo \\
3 & I find it easy to do more than one thing at once. & Agree $\rightarrow$ $\downarrow$ Autismo \\
4 & If there is an interruption, I can switch back to what I was doing very quickly. & Agree $\rightarrow$ $\downarrow$ Autismo \\
5 & I find it easy to read between the lines when someone is talking to me. & Agree $\rightarrow$ $\downarrow$ Autismo \\
6 & I know how to tell if someone listening to me is getting bored. & Agree $\rightarrow$ $\downarrow$ Autismo \\
7 & When I'm reading a story, I find it difficult to work out the characters' intentions. & Agree $\rightarrow$ $\uparrow$ Autismo \\
8 & I like to collect information about categories of things. & Agree $\rightarrow$ $\uparrow$ Autismo \\
9 & I find it easy to work out what someone is thinking or feeling just by looking at their face. & Agree $\rightarrow$ $\downarrow$ Autismo \\
10 & I find it difficult to work out people's intentions. & Agree $\rightarrow$ $\uparrow$ Autismo \\
\bottomrule
\end{longtable}

\subsubsection*{Relación de variables binarias con la salida}
\begin{longtable}{p{0.50\textwidth} p{0.40\textwidth}}
\toprule
Característica & Dirección \\
\midrule
Jaundice (ictericia) & yes $\rightarrow$ $\uparrow$ Autismo \\
Family pdd & yes $\rightarrow$ $\uparrow$ Autismo \\
Used app before & yes $\rightarrow$ $\uparrow$ Autismo \\
Gender & m $\rightarrow$ $\downarrow$ Autismo \\
\bottomrule
\end{longtable}

\subsubsection*{Relación de las variables restantes}
Nota: En etnia, país de residencia y relación hay pocos datos por categoría,
por lo que las tendencias pueden ser engañosas.
\begin{itemize}
    \item \textbf{Ethnicity}: $\uparrow$ Autismo en \emph{black}, \emph{hispanic}, \emph{latino}, \emph{white-european};
    $\downarrow$ Autismo en \emph{middleeastem}, \emph{pasifika}, \emph{southasian}, \emph{turkish}.
    \item \textbf{Age group}: $\uparrow$ Autismo a partir de 30; $\downarrow$ Autismo por debajo de 30.
    \item \textbf{Country of residence} (top 10): $\uparrow$ Autismo en \emph{australia}, \emph{canada},
    \emph{unitedkingdom}, \emph{unitedstates}; $\downarrow$ Autismo en \emph{afghanistan}, \emph{india},
    \emph{jordan}, \emph{newzealand}, \emph{srilanka}, \emph{unitedarabemirates}.
    \item \textbf{Relation}: no aporta información suficiente, por lo que se elimina en etapas posteriores.
\end{itemize}

\section{Rendimiento y validación}
\subsection{Métricas de rendimiento}
Se consideran métricas dependientes e independientes del umbral:
\begin{itemize}
    \item \textbf{Dependientes del umbral}: Accuracy y F1-score.
    \item \textbf{Independientes del umbral}: AUC-ROC y AUC de Precisión-Recall.
\end{itemize}

Se prioriza el \textbf{F1-score} frente a Accuracy debido al desbalance de clases, ya que el F1
combina precisión y recall en la clase positiva. Para métricas independientes del umbral se
elige la \textbf{AUC de Precisión-Recall} en lugar de la ROC, porque es coherente con la
prioridad de detectar bien la clase positiva y es menos optimista en escenarios desbalanceados
(desbalance aproximado $\approx$ 27/73).

\subsection{Metodología utilizada en validación}
Se adopta un enfoque de \textbf{Pipeline + GridSearchCV}. El dataset tiene alrededor de
\textbf{700 instancias}, lo que permite realizar validación cruzada sin un coste computacional
excesivo.

\textbf{Razones para usar Pipeline:}
\begin{itemize}
    \item Encadena preprocesado y modelado en un solo objeto.
    \item Evita fuga de información al ajustar cada transformación solo con datos de entrenamiento.
    \item Hace el flujo reproducible y mantenible.
\end{itemize}

\textbf{Razones para usar GridSearchCV:}
\begin{itemize}
    \item Explora sistemáticamente combinaciones de hiperparámetros.
    \item Es viable con validación cruzada de 10 folds en este tamaño de datos.
    \item Devuelve el mejor modelo reentrenado y métricas por combinación para comparar decisiones.
\end{itemize}

Se consideran distintos esquemas de validación (Hold-Out, K-Fold, Leave-One-Out), y se selecciona
\textbf{StratifiedKFold} para mantener la proporción de clases en cada fold.

\section{Experimentos}
\subsection{Descripción del pipeline}
Se construye la infraestructura para resolver el problema y se definen dos conjuntos de datos:
\textbf{test} y \textbf{resto}. La partición se realiza con una proporción 80/20 y estratificación
por clase, codificando la variable objetivo a valores binarios.

\subsection{Componentes}
Los componentes principales del flujo son:
\begin{itemize}
    \item \textbf{Imputer}: trata valores perdidos en variables numéricas y categóricas.
    \item \textbf{Transformador}: convierte variables categóricas a representación numérica
    (métodos: orden, conteo, \emph{one-hot} y binario).
    \item \textbf{OutlierDetector}: detecta y reemplaza valores atípicos.
    \item \textbf{Estandarizador}: escala variables numéricas para mejorar el comportamiento de KNN.
    \item \textbf{Modelo}: clasificador base, con dos alternativas: KNN o Árbol de decisión.
\end{itemize}

Para el modelo se consideran los hiperparámetros más relevantes:
\begin{itemize}
    \item \textbf{KNN}: \texttt{n\_neighbors}, \texttt{weights} (uniform, distance) y
    \texttt{metric} (manhattan, euclidean).
    \item \textbf{Árbol de decisión}: \texttt{criterion} (gini, entropy), \texttt{min\_samples\_split},
    \texttt{min\_samples\_leaf}, \texttt{max\_depth}, \texttt{ccp\_alpha} y \texttt{random\_state}.
\end{itemize}

\subsection{Construcción de prototipos}
Se define una \emph{pipeline} base con \textbf{Imputer + Transformador}. A partir de ella se construyen
prototipos ampliados incorporando opcionalmente \textbf{OutlierDetector} y \textbf{Estandarizador}, y se
elige un clasificador \textbf{KNN} o \textbf{Árbol de decisión}. Para cada componente se fija un espacio
de búsqueda de hiperparámetros.

\begin{longtable}{p{0.25\textwidth} p{0.68\textwidth}}
\toprule
Componente & Espacio de búsqueda \\
\midrule
Imputer & \texttt{metodo\_imputacion\_vars\_num}: media, mediana.\\
& \texttt{metodo\_imputacion\_vars\_cat}: moda, missing. \\
\addlinespace
Transformador & \texttt{metodo\_cat\_num}: orden, conteo, ohe, binary. \\
\addlinespace
OutlierDetector & $k \in \{0.5, 1.0, \ldots, 9.5\}$; \\ & \texttt{deteccion}: iqr, mediastd; \\ & \texttt{reemplazo}: mediana, media, min, max, moda. \\
\addlinespace
Estandarizador & \texttt{metodo}: minmax, standard, robust. \\
\addlinespace
KNN & \texttt{n\_neighbors}: 3, 5, 7, 9, 11, 13; \\ & \texttt{weights}: uniform, distance; \\ & \texttt{metric}: manhattan, euclidean. \\
\addlinespace
Árbol de decisión & \texttt{criterion}: gini, entropy; \\ & \texttt{min\_samples\_split}: 0.01, 0.05, 0.1; \\ & \texttt{min\_samples\_leaf}: 0.01, 0.05, 0.1; \\ & \texttt{max\_depth}: 2, 4, \ldots, 18; \\ & \texttt{ccp\_alpha}: 0, 0.1, 0.2. \\
\bottomrule
\end{longtable}

\subsection{Análisis de la pipeline base}
Contenido pendiente de completar.

\subsection{Análisis del componente OutlierDetector}
Contenido pendiente de completar.

\subsection{Análisis del componente Estandarizador}
Contenido pendiente de completar.

\subsection{Análisis del componente Modelo}
Contenido pendiente de completar.

\subsection{Análisis del umbral}
Contenido pendiente de completar.

\subsection{Mejor pipeline}
Contenido pendiente de completar.

\section{Resultados}
Contenido pendiente de completar.

\section{Conclusiones}
Contenido pendiente de completar.

\clearpage
\bibliographystyle{plain}
\bibliography{references}

\end{document}
